\documentclass[10pt,a4paper]{amsart}
\usepackage[margin=19mm]{geometry}
\usepackage{amsmath,amssymb,mathtools}
\usepackage[scr=rsfs,cal=euler]{mathalfa}
\usepackage{xfrac}
\usepackage[unicode,hidelinks,bookmarks=false]{hyperref}
\newcommand{\ie}{i.e.\ }
\newcommand{\cf}{cf.\ }
\newcommand{\resp}{respectively\ }
\newcommand{\Subsection}{\subsection}
\providecommand{\nobreakdash}{\nobreak\hbox{-}}
\setlength{\emergencystretch}{2em}
\multlinegap0pt
\usepackage{bm}
\usepackage{bbm}
\usepackage{dsfont}
\usepackage{xspace}
\usepackage{cancel}
\usepackage{mathtools}
\usepackage{amsthm}
\makeatletter
\@ifundefined{theorem}{\newtheorem{theorem}{Theorem}}{}
\@ifundefined{lemma}{\newtheorem{lemma}[theorem]{Lemma}}{}
\makeatother
\newcommand{\PP}{\mathbb P}
\title[The cycle class of the supersingular locus of principally po]{The cycle class of the supersingular locus of principally polarized abelian varieties}
\author{Gerard van der Geer, Shushi Harashita}
\date{Source: arXiv:2307.14393v2 (CC BY 4.0)}
\begin{document}
\maketitle

\noindent\emph{Source and licence.} The cycle class of the supersingular locus of principally polarized abelian varieties. Version arXiv:2307.14393v2 (CC BY 4.0), distributed under the Creative Commons Attribution 4.0 International licence, as recorded in the arXiv licence metadata for this exact version and shown on the abstract page https://arxiv.org/abs/2307.14393v2. Full rights notice: licenses/arxiv-2307.14393-CC-BY-4.0.txt.\par\smallskip

\noindent\textit{Editorial scope.} Counted unit: the Lemma whose source label is \texttt{None} in the article (source0.tex lines 1632--1638, source label \texttt{None}), delivered together with the article's own complete original proof of it (source0.tex lines 1639--1678, one \texttt{proof} environment of 40 lines). No mathematics is written, repaired or re-derived: statement and proof are the article's own text line for line. Required context retained uncounted: none. Every definition, display and label of those lines is retained unchanged. Omitted from this package and not redistributed: 1--1631, 1679--2753. Nothing inside a retained excerpt is omitted or altered: every retained body is delivered exactly as the article's lines are written, so no omission receipt of any kind accompanies this package. Citations: the retained excerpts carry no citation command at all, so no bibliography entry of the article is transcribed and no reference is redistributed. Reference adaptation: none was needed and none was made; every reference inside the delivered unit resolves inside it, so no unresolved reference or citation is produced. Printed slips kept literal: no printed word, symbol, hypothesis or number of the counted statement or of its proof was altered. Layout and class: the article's own class is \texttt{amsart} with packages including amssymb, amscd, latexsym, epsfig, extarrows, hyperref, xy, pst-all; the wrapper is the lane's pinned \texttt{amsart} editorial wrapper at 10pt on A4, which restates the theorem-like environments the retained text needs and adds the article's own macro definitions that the retained text needs (\texttt{\textbackslash PP}), with extra wrapper packages (bm, bbm, dsfont, xspace, cancel, mathtools). The delivered numbering is the unit's own. Assets: no retained span contains a figure environment, a graphics-inclusion command, a PDF-page inclusion, a table or a TikZ picture; no image file is copied into this package, the package declares no asset dependency and no image file is redistributed with it. Licence: version arXiv:2307.14393v2 of this article is distributed under the Creative Commons Attribution 4.0 International licence, as recorded in the arXiv licence metadata for this exact version and shown on the abstract page https://arxiv.org/abs/2307.14393v2. A full rights notice is delivered at licenses/arxiv-2307.14393-CC-BY-4.0.txt. Characters: every retained excerpt is pure ASCII, so the delivered engine, XeLaTeX with its default Latin Modern text font, reports no missing character.
\begin{lemma}
The threefold $\mathcal{F}_1$ is a divisor in a ${\PP}^2$-bundle ${\PP}(B)$
with $B$ the rank $3$ vector bundle defined by $M_2/\tilde{L}$
over a surface $\tilde{\mathcal{F}}_2^{\prime}$ obtained by an inseparable base change 
$\tilde{\mathcal{F}}_2^{\prime} \to\tilde{\mathcal{F}}_2$
of degree $p$.
\end{lemma}

\begin{proof}
Recall that in order to define $M_1 \subset M_2$, we chose a basis 
$$
v=a_5v_0+a_6Fx_2+a_7Fx_3+a_8Fx_4, \quad
w=a_9Fx_2+a_{10}Fx_3+a_{11}Fx_4 \, 
$$
as in (6) with $\langle v,Fv\rangle$, $\langle v,Fw\rangle$, $\langle Fv,w\rangle$
all in $W$. 
The equations ($g_2$) and ($g_3$) correspond to the inseparable 
base change $\tilde{\mathcal{F}}_2^{\prime} \to\tilde{\mathcal{F}}_2$
given on the locus with $a_1\neq 0$ by (11) 
$$
(a_9/a_{10})^p= (a_2-a_2^{p^2})/(a_3-a_3^{p^2})\, .
$$
Then on $\tilde{\mathcal{F}}_2^{\prime}$ we have the bundle ${\PP}(B)$.
If $a_5\neq 0$ the morphism
${\mathcal F}_1 \to \tilde{\mathcal F}_2$ is defined by sending
$(M_1 \subset M_2 \subset M_3)$ to the point defined by $L:=M_1 \cap V^{-1}M_2^t \bmod FM_2$.
Indeed,
by $(F,V)M_2 \subset M_1$, the subspace $L$ contains $(F,V)M_2 \bmod FM_2$,
and $L$ is the one-dimensional space generated by $w$ of (6),
since one can check $\langle Vw, M_2 \rangle\subset W$ and if $a_5 \ne0$, 
then $\langle Vv, M_2 \rangle \not\subset W$.

For $a_5\neq 0$ we find from $\langle v, Fv\rangle \in W$ an equation 
$$
a_1a_8^p+a_2a_7^p-a_3a_6^p -a_5^{p-1}(a_1^pa_8+a_2^pa_7-a_3^pa_6)=0\, .
$$
This defines a rational curve with a cusp in ${\PP}^2={\PP}(M_2/\tilde{L})$.
As $\mathcal{F}_1$ is defined as the closure of the space of rigid flags,
we obtain that this equation defines $\mathcal{F}_1$ in ${\PP}(B)$.
Observe that in order to analyze this we may assume as we did in the 
preceding section that
$a_1=1$ and $a_2\neq 0$ and then $a_6=0$ and the curve can be written 
in coordinates $(a_5:a_7:a_8)$ as
$$
a_1a_8^p+a_2a_7^p -a_5^{p-1}(a_1^pa_8+a_2^pa_7)=0\, .
$$
The cusp is determined by $a_5=0$ and $a_8+a_2^{1/p}a_7=0$. 
\end{proof}

\end{document}
